% ============================================================
%  informe.tex — Plantilla de informe técnico UCuenca
%  Módulo 1217 · Redes Complejas
%  Estilo basado en el PDF de referencia del proyecto
% ============================================================

\documentclass[12pt, a4paper]{article}

% ============================================================
% Codificación y lenguaje
% ============================================================
\usepackage[T1]{fontenc}
\usepackage[utf8]{inputenc}
\usepackage[spanish, es-nodecimaldot]{babel}

% ============================================================
% Fuente — Palatino (cuerpo + matemáticas)
% ============================================================
\usepackage{mathpazo}
\usepackage[scaled=0.92]{helvet}   % Helvetica para código inline si se desea
\usepackage{microtype}             % Mejora el espaciado tipográfico

% ============================================================
% Geometría de página
% ============================================================
\usepackage[
    a4paper,
    left   = 2.6cm,
    right  = 2.6cm,
    top    = 3.0cm,
    bottom = 3.0cm
]{geometry}

% ============================================================
% Matemáticas
% ============================================================
\usepackage{amsmath}
\usepackage{amssymb}
\usepackage{bm}            % Negritas en matemáticas

% ============================================================
% Gráficos y figuras
% ============================================================
\usepackage{graphicx}
\usepackage{float}
\usepackage{subcaption}    % Para subfiguras (a), (b)

% ============================================================
% Colores
% ============================================================
\usepackage{xcolor}

%  Azul principal de secciones (igual al del PDF de referencia)
\definecolor{azulsec}{RGB}{26, 88, 140}

%  Fondo y borde del resumen
\definecolor{grisfondo}{RGB}{245, 245, 245}
\definecolor{borde}{RGB}{70, 70, 70}

%  Encabezado de tablas
\definecolor{tabhead}{RGB}{30, 70, 115}      % Azul oscuro → letras blancas
\definecolor{tabfila}{RGB}{235, 241, 248}    % Azul muy claro → filas alternas

%  Caja de aviso (naranja, como los boxes del PDF)
\definecolor{naranja}{RGB}{230, 150, 0}
\definecolor{naranjaFondo}{RGB}{255, 243, 210}

% ============================================================
% Tablas
% ============================================================
\usepackage{booktabs}      % \toprule, \midrule, \bottomrule
\usepackage{colortbl}      % Colorear filas/columnas
\usepackage{array}         % Tipos de columna extendidos
\usepackage{tabularx}      % Columnas de ancho automático
\usepackage{multirow}      % Celdas multifila
\usepackage{makecell}      % Saltos de línea en celdas

\renewcommand{\arraystretch}{1.35}   % Altura de filas

% Tipo de columna centrado con ancho fijo
\newcolumntype{C}[1]{>{\centering\arraybackslash}p{#1}}
% Tipo de columna justificado con ancho fijo
\newcolumntype{L}[1]{>{\raggedright\arraybackslash}p{#1}}
% Tipo de columna derecha con ancho fijo
\newcolumntype{R}[1]{>{\raggedleft\arraybackslash}p{#1}}

% Comando para la fila de encabezado de tabla:
% Uso: \encabezado{col1} & \encabezado{col2} & ... \\
\newcommand{\encabezado}[1]{%
    \cellcolor{tabhead}\textcolor{white}{\textbf{#1}}}

% ============================================================
% Captions (pies de figura y tabla)
% ============================================================
\usepackage{caption}
\captionsetup{
    font        = small,
    labelfont   = bf,
    labelsep    = period,
    justification = justified,
    singlelinecheck = false
}
\captionsetup[figure]{
    name       = Fig.,
    justification = centering
}
\captionsetup[table]{
    name       = Tabla,
    position   = above,
    justification = centering
}

% ============================================================
% Formato de secciones — titlesec
% ============================================================
\usepackage{titlesec}

%  Sección: número y título en azul, negrita, subrayado con línea
\titleformat{\section}
    {\color{azulsec}\large\bfseries}   % Formato del texto
    {\thesection.}                      % Número + punto
    {0.5em}                             % Separación número-título
    {}                                  % Código antes del título
    [\color{azulsec}\vspace{2pt}\rule{\linewidth}{0.7pt}]  % Regla debajo

%  Subsección: número y título en azul, negrita, sin regla
\titleformat{\subsection}
    {\color{azulsec}\normalsize\bfseries}
    {\thesubsection}
    {0.5em}
    {}

%  Subsubsección: azul, normalsize, negrita itálica
\titleformat{\subsubsection}
    {\color{azulsec}\normalsize\bfseries\itshape}
    {\thesubsubsection}
    {0.5em}
    {}

%  Sección sin número (para Anexos, Referencias manuales, etc.)
%  Uso: \section*{Anexos}
\titleformat{name=\section, numberless}
    {\color{azulsec}\large\bfseries}
    {}
    {0em}
    {}
    [\color{azulsec}\vspace{2pt}\rule{\linewidth}{0.7pt}]

%  Espaciado antes/después de secciones
\titlespacing{\section}    {0pt}{16pt}{6pt}
\titlespacing{\subsection} {0pt}{10pt}{4pt}
\titlespacing{\subsubsection}{0pt}{8pt}{3pt}

% ============================================================
% Entorno de Resumen con borde izquierdo y fondo gris
% ============================================================
\usepackage[framemethod=default]{mdframed}

\newenvironment{resumen}{%
    \begin{mdframed}[
        backgroundcolor = grisfondo,
        linecolor       = borde,
        linewidth       = 0pt,
        leftline        = true,
        rightline       = false,
        topline         = false,
        bottomline      = false,
        innerleftmargin = 10pt,
        innerrightmargin= 10pt,
        innertopmargin  = 8pt,
        innerbottommargin = 8pt,
        leftmargin      = 0pt,
        rightmargin     = 0pt,
        skipabove       = 10pt,
        skipbelow       = 10pt,
        leftmarginwidth = 4pt     % Grosor del borde izquierdo
    ]
    \small
}{%
    \end{mdframed}
}

% ============================================================
% Caja de aviso / nota (borde naranja)
% ============================================================
\newenvironment{aviso}{%
    \begin{mdframed}[
        backgroundcolor = naranjaFondo,
        linecolor       = naranja,
        linewidth       = 1pt,
        leftline        = true,
        rightline       = true,
        topline         = true,
        bottomline      = true,
        innerleftmargin = 8pt,
        innerrightmargin= 8pt,
        innertopmargin  = 5pt,
        innerbottommargin = 5pt,
        skipabove       = 6pt,
        skipbelow       = 6pt
    ]
    \small
}{%
    \end{mdframed}
}

% ============================================================
% Código (monoespaciado inline y en bloque)
% ============================================================
\usepackage{listings}
\usepackage{upquote}

\lstset{
    basicstyle      = \small\ttfamily,
    keywordstyle    = \color{azulsec}\bfseries,
    commentstyle    = \color{gray}\itshape,
    stringstyle     = \color{teal},
    backgroundcolor = \color{grisfondo},
    frame           = single,
    framerule       = 0.4pt,
    rulecolor       = \color{gray!40},
    breaklines      = true,
    breakatwhitespace = true,
    showstringspaces = false,
    numbers         = left,
    numberstyle     = \tiny\color{gray},
    numbersep       = 5pt,
    tabsize         = 4,
    captionpos      = b
}

% Código inline rápido: \code{texto}
\newcommand{\code}[1]{\texttt{\small #1}}

% ============================================================
% Listas
% ============================================================
\usepackage{enumitem}
\setlist[itemize]{
    leftmargin = 1.5em,
    itemsep    = 2pt,
    topsep     = 4pt
}
\setlist[enumerate]{
    leftmargin = 1.5em,
    itemsep    = 2pt,
    topsep     = 4pt
}

% ============================================================
% Párrafos — sin sangría, con espacio entre párrafos
% ============================================================
\setlength{\parindent}{0pt}
\setlength{\parskip}{6pt}

% ============================================================
% Encabezado y pie de página
% ============================================================
\usepackage{fancyhdr}
\pagestyle{fancy}
\fancyhf{}
\renewcommand{\headrulewidth}{0pt}
\fancyfoot[C]{\small\thepage}

% ============================================================
% Hipervínculos (cargar siempre al final)
% ============================================================
\usepackage{hyperref}
\hypersetup{
    colorlinks  = true,
    linkcolor   = azulsec,
    citecolor   = azulsec,
    urlcolor    = azulsec,
    pdftitle    = {Informe Técnico},
    pdfauthor   = {Universidad de Cuenca},
    pdfsubject  = {Redes Complejas},
    pdfkeywords = {redes complejas, UCuenca}
}

% ============================================================
% Comandos auxiliares
% ============================================================

% Separador de autores (centrado, con punto medio)
\newcommand{\autorsep}{\quad\textperiodcentered\quad}

% Negrita en ecuaciones
\newcommand{\vect}[1]{\boldsymbol{#1}}

% Promedio
\newcommand{\prom}[1]{\langle #1 \rangle}

% Norma
\newcommand{\norm}[1]{\left\| #1 \right\|}

% ============================================================
%  ██████████████████████████████████████████████████████████
%  INICIO DEL DOCUMENTO
%  ██████████████████████████████████████████████████████████
% ============================================================
\begin{document}

% ------------------------------------------------------------
% BLOQUE DE TÍTULO
% Reemplaza los campos con los datos del proyecto
% ------------------------------------------------------------
\begin{center}
    % Título principal (LARGE negrita, multilínea si es largo)
    {\LARGE\bfseries
        Título del informe aquí\\[0.3em]
        Segunda línea del título si es necesaria%
    }\\[0.8em]

    % Subtítulo / descripción del proyecto
    {\normalsize
        Subtítulo o descripción del proyecto — Módulo 1217, Opción X%
    }\\[0.8em]

    % Autores separados por punto medio
    {\normalsize
        Autor Uno \autorsep Autor Dos \autorsep Autor Tres%
    }\\[0.3em]

    % Institución, correo y fecha
    {\small
        Universidad de Cuenca
        \enspace---\enspace correo@ucuenca.edu.ec
        \enspace---\enspace mes de 2026%
    }
\end{center}

\vspace{0.4em}

% ------------------------------------------------------------
% RESUMEN
% Comienza con \textbf{Resumen---} según el estilo del PDF
% ------------------------------------------------------------
\begin{resumen}
\textbf{Resumen---} Escribir aquí el resumen del trabajo. Describir brevemente
el contexto del problema, la metodología utilizada, los resultados principales
obtenidos y la conclusión más relevante. Puede incluir valores numéricos
destacados en \textbf{negrita}.
\end{resumen}

\vspace{0.2em}

% ============================================================
% SECCIONES — El usuario añade su contenido aquí
% ============================================================

% ------------------------------------------------------------
\section{Introducción}
% ------------------------------------------------------------

Texto de la introducción. Se puede citar referencias con~\cite{ref1}.
Texto en \textit{cursiva}, en \textbf{negrita} o en \code{monoespaciado}.

Segundo párrafo de la introducción.

% ------------------------------------------------------------
\section{Descripción Metodológica}
% ------------------------------------------------------------

Texto introductorio de la sección.

\subsection{Primera subsección}

Texto de la subsección. Ejemplo de ecuación numerada:

\begin{equation}
    E_{\text{total}}(n_i) = E_{\text{line}} + E_{\text{turn}} + E_{\text{star}}
    \label{eq:energia_total}
\end{equation}

Referencia a la ecuación: véase la ec.~\eqref{eq:energia_total}.

\subsection{Segunda subsección}

Ejemplo de ecuación sin número:

\begin{equation*}
    F(q_i) = G(q_i) + H(q_i)
\end{equation*}

Ejemplo de alineación de ecuaciones:

\begin{align}
    G(q_i) &= G(q_{i-1}) + E_{\text{total}}(n_i) \label{eq:G} \\
    H(q_i) &= \left[\mu_r (M_{\text{car}} + m_{\text{load}}) g
               + \frac{P_{\text{base}}}{v} \right]
               \cdot \norm{p_{\text{goal}} - p_i}_2 \label{eq:H}
\end{align}

% ----  EJEMPLO DE TABLA ----
\begin{table}[H]
\centering
\caption{Parámetros utilizados en el modelo (valores de ejemplo).}
\label{tab:parametros}
\begin{tabular}{lccc}
    \toprule
    \encabezado{Parámetro} &
    \encabezado{Símbolo}   &
    \encabezado{Valor A}   &
    \encabezado{Valor B}   \\
    \midrule
    Primer parámetro  & $P_1$ & 1.0  & 0.5  \\
    Segundo parámetro & $P_2$ & 2.0  & 1.5  \\
    Tercer parámetro  & $P_3$ & 18 kg & 0.28 kg \\
    Cuarto parámetro  & $P_4$ & 0--8 kg & 0--80 g \\
    \bottomrule
\end{tabular}
\end{table}

% ----  EJEMPLO DE TABLA CON FILAS ALTERNAS ----
\begin{table}[H]
\centering
\caption{Resultados comparativos (filas alternas coloreadas).}
\label{tab:resultados}
\begin{tabularx}{\linewidth}{L{3.5cm} C{2.5cm} C{2.5cm} C{2.5cm}}
    \toprule
    \encabezado{Métrica} &
    \encabezado{Método A} &
    \encabezado{Método B} &
    \encabezado{$\Delta$ (\%)} \\
    \midrule
    \rowcolor{tabfila}
    Distancia (m)   & 34.02  & 35.31  & $+3.8\,\%$ \\
    Giro (rad)      & 27.49  & 19.63  & $-28.6\,\%$ \\
    \rowcolor{tabfila}
    Energía (kJ)    & 3.42   & 3.03   & $-11.2\,\%$ \\
    \bottomrule
\end{tabularx}
\end{table}

% ----  EJEMPLO DE FIGURA ----
\begin{figure}[H]
    \centering
    \includegraphics[width=0.7\linewidth]{images/figura1.png}
    \caption{Descripción breve de la figura. Puede ocupar más de una
             línea si la explicación es extensa.}
    \label{fig:figura1}
\end{figure}

% ----  EJEMPLO DE SUBFIGURAS ----
\begin{figure}[H]
    \centering
    \begin{subfigure}[b]{0.47\linewidth}
        \centering
        \includegraphics[width=\linewidth]{images/figura2a.png}
        \caption{Subfigura izquierda.}
        \label{fig:fig2a}
    \end{subfigure}
    \hfill
    \begin{subfigure}[b]{0.47\linewidth}
        \centering
        \includegraphics[width=\linewidth]{images/figura2b.png}
        \caption{Subfigura derecha.}
        \label{fig:fig2b}
    \end{subfigure}
    \caption{Título general de las dos subfiguras.}
    \label{fig:figura2}
\end{figure}

\subsection{Tercera subsección con caja de aviso}

Texto antes del aviso.

\begin{aviso}
\textbf{Nota:} Este es un cuadro de aviso con borde naranja, útil para
señalar elementos pendientes, advertencias o notas importantes.
\end{aviso}

Texto después del aviso.

% ------------------------------------------------------------
\section{Análisis de Resultados}
% ------------------------------------------------------------

\subsection{Primera subsección de resultados}

Texto de análisis de resultados. Se puede hacer referencia a la
Tabla~\ref{tab:resultados} y a la Fig.~\ref{fig:figura1}.

Ejemplo de lista con viñetas:

\begin{itemize}
    \item \textbf{Primer ítem:} descripción del ítem en negrita seguida
          de texto explicativo que puede extenderse en múltiples líneas.
    \item \textbf{Segundo ítem:} descripción correspondiente.
    \item \textbf{Tercer ítem:} descripción correspondiente.
\end{itemize}

Ejemplo de lista numerada:

\begin{enumerate}
    \item Primer paso del procedimiento.
    \item Segundo paso del procedimiento.
    \item Tercer paso del procedimiento.
\end{enumerate}

% ------------------------------------------------------------
\section{Conclusiones}
% ------------------------------------------------------------

\begin{itemize}
    \item Primera conclusión principal obtenida del trabajo.
    \item Segunda conclusión principal.
    \item Tercera conclusión y trabajo futuro.
\end{itemize}

% ------------------------------------------------------------
\section*{Referencias}
% ------------------------------------------------------------

% Estilo numerado entre corchetes [1], [2], etc.
% Agrega {99} si hay más de 9 referencias (para alinear correctamente)
\begin{thebibliography}{9}

\bibitem{ref1}
A. Apellido y B. Apellido,
``Título del artículo de referencia,''
\textit{IEEE Trans. Nombre Revista},
vol.~X, n.\textsuperscript{o}~Y, pp.~Z--ZZ, mes año.

\bibitem{ref2}
C. Apellido \textit{et al.},
``Título del segundo artículo,''
en \textit{Proc. Nombre Conferencia},
Ciudad, País, año, pp.~1--6.

\bibitem{ref3}
D. Apellido,
\textit{Título del Libro}.
Ciudad: Editorial, año.

\end{thebibliography}

% ------------------------------------------------------------
\section*{Anexos}
% ------------------------------------------------------------

\textbf{A. Primer Anexo:} descripción del contenido del primer anexo.
Puede incluir código fuente, tablas adicionales o figuras complementarias.

\vspace{0.5em}
\textbf{B. Código fuente:} repositorio disponible en
\url{https://github.com/usuario/repositorio}.
Instrucciones de instalación y reproducción en \code{INSTALACION.md}.

\vspace{0.5em}
\textbf{C. Resultados adicionales:} archivos \code{*.csv} y \code{*.png}
disponibles en la carpeta \code{results/}.

% ============================================================
\end{document}
% ============================================================
